% !PRINTED 455
\weberpage{455}
%<<CONT
follows, and if one again sets

\weq{29}{\beta c'-\beta c+\alpha p+\delta p=y}

then:

\weq{30}{\begin{aligned} 2(\beta c'+\alpha p) &= Bx+y,\\ 2(\beta c-\delta p) &= Bx-y. \end{aligned}}

From this

\weq{31}{y^2-Dx^2=4p^2.}

From (28) there result the congruences

\weq{32}{\begin{aligned} \beta c c' &\equiv Cx, \quad \beta(c+c') \equiv Bx, \quad \beta \equiv Ax\\ Dx^2 &\equiv \beta^2(c-c')^2 \pmod{p}. \end{aligned}}

If $c'$ is to be different from $c$ modulo $p$, then according to this neither $x$ nor $D$ can be divisible by $p$, and it follows from (32)

\weq{33}{\begin{aligned} A(c+c') &\equiv B,\\ Acc' &\equiv C, \end{aligned}}

and this leads, by elimination of $c'$, to the congruence (13) for $c$, and so to the value of $c$ that is to be excluded.

\gesperrt{If we summarise this, we see that among the quantities (15) equivalent ones occur only in the two exceptional cases $D=-3$, $D=-4$, and indeed in the case $D=-3$ three of them at a time, in the case $D=-4$ two of them at a time, are equivalent.}

In these cases, therefore, the right-hand side of formula (16) must still be divided by 3 or by 2, and the formula thus remains correct even then if, as already earlier, for $D=-3, -4$ one understands by $h$ not 1, but $\frac{1}{3}$ and $\frac{1}{2}$.

\begin{center}
\textbf{\S\ 124. \gesperrt{Class Fields and Order Fields.}}
\end{center}

If $D=Q^2\Delta$ is a discriminant and $\Delta$ its stem, then by the results of the preceding section the class invariant $u$ of the discriminant $\Delta$ can be expressed rationally by the class invariant $v$ of the discriminant $D$, and to every $u$ there belong $r$ values of $v$, if $r$ is determined by formula (16), \S\,123. These values of $v$ are the roots of an equation of degree $r$, whose coefficients depend rationally on $u$. The quantities $v$ therefore belong to a field $\mathfrak{K}(D)$ over $\mathfrak{K}(\Delta)$, which is none other than the class field of the discriminant $D$.

With reference to the quadratic fields $\Omega=\mathfrak{R}(\sqrt{\Delta})$ both fields are \gesperrt{Abelian}. We wish, if a more exact investi-
%CONT>>
